Generating synthetic financial time series that faithfully preserve the statistical properties of real market data is a central challenge in quantitative finance.
High-fidelity synthetic data are essential for stress testing risk models against plausible but unobserved market scenarios, validating portfolio optimization algorithms, and augmenting limited training sets for machine learning pipelines \citep{assefa2020generating, jordon2022synthetic}.
Empirical equity excess growth rates exhibit three well-documented statistical regularities---heavy-tailed (leptokurtic) distributions, negligible linear autocorrelation in raw returns, and persistent volatility clustering---that any generative framework must simultaneously reproduce to be considered statistically faithful \citep{cont2001empirical, mandelbrot1963variation}.

Parametric approaches to synthetic data generation face a fundamental tension between distributional accuracy and temporal fidelity.
GARCH models \citep{engle1982autoregressive, bollerslev1986generalized} excel at reproducing the slow decay of the autocorrelation function (ACF) of absolute returns---the hallmark of volatility clustering---but their unimodal innovation distribution cannot represent the multi-regime structure of equity markets, leading to poor distributional fit.
Conversely, i.i.d.\ generators that match the marginal distribution (e.g., Laplace sampling) achieve high distributional pass rates but produce flat ACF profiles by construction.
Hidden Markov models (HMMs) offer a middle path: the latent regime structure can simultaneously generate heavy-tailed mixtures and temporal persistence through the transition dynamics.

However, the application of HMMs to financial stylized facts has been constrained by an influential negative result.
\citet{ryden1998stylized} systematically evaluated Gaussian-emission HMMs with $K = 2$--$3$ states against the canonical stylized facts and concluded that while HMMs reproduce heavy tails and negligible return autocorrelation, they \emph{fail} to capture the slow decay of the ACF of absolute returns.
The mechanism is clear: with few states, the geometric sojourn-time distribution of a standard Markov chain produces exponentially fast ACF decay, insufficient to match the hyperbolic-like persistence observed empirically.
\citet{bulla2006stylized} subsequently showed that hidden semi-Markov models (HSMMs), which replace the geometric sojourn distribution with flexible alternatives, restore ACF fidelity---but at the cost of substantially increased model complexity.

\citet{alswaidan2026hybrid} proposed an alternative solution within the discrete HMM framework: a Poisson jump-duration mechanism that forces the model to dwell in tail states for empirically calibrated durations.
That approach achieved 91--95\% Kolmogorov--Smirnov (KS) pass rates across 424 US equities and partially reproduced volatility clustering.
However, the discrete framework introduced three sources of complexity: (i)~Laplace quantile binning, which discretizes continuous observations and introduces quantization error; (ii)~two additional hyperparameters ($\epsilon$ for jump probability and $\lambda$ for mean jump duration) requiring grid search calibration; and (iii)~within-bin Student-$t$ emission models whose parameters are estimated separately from the transition structure.

In this paper, we show that these complexities are unnecessary.
A continuous Gaussian HMM (CHMM) trained via the Baum-Welch algorithm \citep{baum1970maximization, rabiner1986introduction} with quantile-based initialization reproduces all three canonical stylized facts at moderate state resolution ($K = 6$--$13$) \emph{without} requiring jump augmentation, discretization, or separate emission fitting.
The key insight is that the negative result of \citet{ryden1998stylized} is an artifact of insufficient state resolution: at $K \geq 3$ with proper initialization, the Baum-Welch algorithm learns a transition matrix with sufficient diagonal dominance in high-volatility states to generate slow ACF decay, effectively approximating the flexible sojourn-time distributions of HSMMs through a mixture of geometric distributions.

The contributions of this work are:
\begin{enumerate}
    \item \textbf{Baum-Welch training without jump augmentation.}
    We train continuous Gaussian emission parameters $(\mu_k, \sigma_k)$ and the transition matrix $\mathbf{T}$ jointly from raw observations via the Baum-Welch algorithm with quantile-based initialization.
    The resulting CHMM achieves 90--94\% KS pass rates and ACF-MAE of 0.046--0.049 on absolute returns---comparable to GARCH(1,1) (ACF-MAE $= 0.047$)---without any jump mechanism, eliminating two hyperparameters ($\epsilon$, $\lambda$) and the discretization step.

    \item \textbf{Contradiction of the Ryd\'{e}n et al.\ (1998) consensus.}
    We demonstrate that the widely cited limitation of Gaussian HMMs---inability to reproduce volatility clustering---is a consequence of testing only $K = 2$--$3$ states.
    At moderate $K$, the learned transition matrix develops sufficient persistence in high-volatility states to achieve ACF-MAE comparable to GARCH, which we show is effectively an approximation to the HSMM solution of \citet{bulla2006stylized}.

    \item \textbf{Comprehensive five-model benchmark.}
    We evaluate the CHMM against bootstrap resampling, Gaussian and Laplace i.i.d.\ generators, and GARCH(1,1) using seven complementary metrics---KS and Anderson--Darling (AD) pass rates, excess kurtosis, ACF-MAE, Wasserstein-1 and Hellinger distances, and quantile coverage---on 1,000 simulated paths in both in-sample (2014--2024) and out-of-sample (2025) windows.

    \item \textbf{State resolution analysis.}
    We systematically vary $K \in \{3, 6, 9, 11, 13\}$ and show that distributional fidelity improves with $K$ while ACF-MAE remains remarkably stable, and quantile coverage is 100\% at all resolutions.

    \item \textbf{Cross-asset and volatility index generalization.}
    We confirm robustness across NVDA, JNJ, and JPM (97--100\% IS KS), and extend the framework to CBOE VIX and VXX log returns, demonstrating that the CHMM captures the distinct statistical properties of volatility indices and provides a natural bridge to regime-switching option pricing.
\end{enumerate}

The remainder of this paper is organized as follows.
Section~\ref{sec:related} reviews the relevant literature on stylized facts, parametric models, HMMs, and deep generative approaches.
Section~\ref{sec:methods} describes the data, the CHMM specification, the Baum-Welch training procedure, the benchmark models, and the validation metrics.
Section~\ref{sec:results} presents the empirical results: descriptive statistics, five-model comparison, state resolution sensitivity, out-of-sample evaluation, cross-asset generalization, and the VIX/VXX extension.
Section~\ref{sec:discussion} interprets the findings---particularly the ACF-MAE result and its implications for the Ryd\'{e}n et al.\ consensus---and identifies limitations.
Section~\ref{sec:conclusion} summarizes the contributions and outlines future directions.
